\documentclass{ximera}

\title{A Moving Image}

\author{Bart Snapp}

\begin{document}
\begin{abstract}
  All translations can be viewed as compositions of reflections.
\end{abstract}
\maketitle

In this activity, we are going to explore how reflections are related
to translations.  Consider the following ``F'' shape:
\[
\includegraphics{planeF.pdf}
\]

\begin{problem}
Use two reflections to translate the ``F'' shape 8 units to the
right. Don't worry about the actual matrices involved.
\end{problem}

\begin{problem}
What reflections would translate the original ``F'' shape 4 units to
the right? Again, don't worry about the actual matrices involved.
\end{problem}

\begin{problem}
Can you generalize the two problems above? What composition of two
reflections will give a horizontal translation of $n$ units to the
right?
\end{problem}

\begin{problem} 
Now suppose you wish to move the original ``F'' shape 4 units to the
right, but you will first reflect via this matrix:
\[
\mat{F}_{x=0}  = 
\begin{bmatrix}
-1 & 0 & 0 \\
0 & 1 & 0 \\
0 & 0 & 1
\end{bmatrix}
\]
What is the other matrix that is needed?
\end{problem}

\begin{problem}
Generalize the problemlem above: suppose you wish to translate the
original ``F'' shape $n$ units to the right, and you start by
reflecting by $\mat{F}_{x=0}$. What is the second matrix that you
apply?
\[
\includegraphics{planeF.pdf}
\]
\end{problem}


\end{document}
